\section{Performance, Scalability, and Critical Analysis}
\label{sec:analysis}

\subsection{What has and has not been benchmarked}

The current implementation has strong correctness instrumentation but limited
performance instrumentation. The fresh test run establishes a 537-test, 166.91-second
verification result under coverage on one development environment. It does not isolate
physics step throughput, renderer cost, UI transport latency, memory, or scaling with
module count. Likewise, the documented approximately 177-s simulated completion is a
scenario outcome, not wall-clock performance.

This distinction is important because the demonstration's purpose is presently
functional: prove that the semantic architecture can support wheel/contact dynamics and
reconfiguration. A journal-quality performance claim will require a controlled
benchmark harness, machine disclosure, repeated trials, and raw traces.

\subsection{Algorithmic cost by subsystem}

Let $M$ be module count, $C$ connector count, $E$ active connection count, $J$ joint
count, and $K$ the number of nearby broad-phase pairs.

\begin{table}[htbp]
\centering
\caption{Expected algorithmic costs of principal semantic operations.}
\label{tab:complexity}
\begin{tabularx}{\textwidth}{L{0.28\textwidth}L{0.22\textwidth}Y}
\toprule
Operation & Expected cost & Comment \\
\midrule
Snapshot ingestion & $O(M+J+C)$ & Copies observed link/joint/frame state \\
Spatial hash build/query & expected $O(C+K)$ & Avoids all-pairs $O(C^2)$ in sparse scenes \\
Acceptance/guards per pair & $O(1)$ & Fixed-size vector/quaternion computations \\
Dock commit & component-set dependent & Backend allocation plus assembly merge \\
Undock/split & $O(|V_A|+|E_A|)$ & Recomputes only affected assembly component \\
Topology view build & $O(M+E)$ & Includes isolated nodes and parallel edges \\
Metric summary & $O(M+C+E+|\mathcal L|)$ in simple aggregation & Event history counters can later be incremental \\
Event delta transport & $O(\Delta |\mathcal L|)$ & Sends only unseen canonical events \\
\bottomrule
\end{tabularx}
\end{table}

The broad phase is the primary semantic scaling mechanism. It is effective when
connectors are spatially sparse relative to capture tolerance. Dense piles or very
large tolerance values increase $K$ and approach all-pairs behavior. This has not yet
been measured empirically.

\subsection{Model-view cache behavior}

The model-view factory retains a bounded least-recently-used set of latest snapshots,
defaulting to 32 entries. A topology view token depends on backend sample and topology
revision, so pose motion rebuilds node values while a pure event revision can reuse
graph content. This prevents unbounded per-step graph accumulation.

The cache is correctness tested, but hit rate, build latency, and serialization cost
have not been benchmarked. At seven modules they are negligible relative to physics;
at thousands of modules or multiple simultaneous views they may become visible.

\subsection{Physics and rendering decoupling}

The physical snake default has $210/0.002=105{,}000$ nominal solver steps. A native
viewer synchronization on every step would couple this to 500 attempted visual updates
per simulated second. Commit \code{4a5ca00} caps native synchronization at about 30 Hz,
while retaining every physics step and forcing a final sync. Runtime Inspector graph
publication defaults to 20 Hz wall time, and event deltas remain lossless between
frames.

This is an architectural improvement for accelerated presentation. It reduces
render-call frequency but does not prove a particular achieved real-time factor. The
solver, scene complexity, contacts, Python controller, IPC, native viewer, and host GPU
all contribute. If they take longer than the requested pacing deadline, the runtime
does not skip physics; it simply falls behind the target.

\subsection{Why playback acceleration preserves the experiment}

Changing $\rho$ in \cref{eq:pacing} changes wall waiting only. Contrast three possible
implementations:

\begin{enumerate}
  \item Increasing solver $\Delta t$ reduces step count and changes contact stability.
  \item Multiplying wheel targets changes robot dynamics and event times.
  \item Dividing wall deadlines by $\rho$ preserves simulated trajectories subject to
        deterministic solver execution.
\end{enumerate}

\modsim{} uses the third. Unit tests validate factors 1, 2, 4, and slow motion, reject
zero/non-finite factors, verify protocol round-trip, and inspect viewer deadlines with a
fake clock. A historical protocol record without the new field defaults to 1.0.

\subsection{Physical model debugging lessons}

The physical scenario exposed issues that kinematic staging cannot reveal.

\paragraph{Visual meshes are not contact models.}
Detailed Fusion STLs are suitable for appearance but expensive and fragile for contact.
The pack now uses visual-only STLs and explicit wheel/skid/face primitives.

\paragraph{Rotating square face proxies obstruct locomotion.}
Early LEFT/RIGHT connector proxies were broad square boxes on rotating wheel bodies.
Their corners swept through neighboring module geometry and stalled connected trains.
Coaxial disk/cylinder proxies preserve face contact without creating paddle-like
interference.

\paragraph{Friction direction must follow the wheel.}
Anisotropic coefficients alone are insufficient if the contact tangent stays in a
world-fixed direction. The backend rotates the tangent basis to follow the wheel axle
at every contact solve.

\paragraph{A weld can overconstrain contact.}
Leaving a direct body-pair contact active while also applying a rigid weld causes two
constraints to fight. Pairing each runtime weld with a preallocated exclusion removes
that specific conflict while leaving all other module and assembly contacts active.

\paragraph{Multi-module differential drive is state dependent.}
A controller that moves one module does not automatically move a rigid three-module
train. Assembly allocation, traction, pivot friction, route hysteresis, and lateral
prebias were required. This reinforces the separation between generic semantic plan
and platform-specific physical route.

\subsection{Stability and safety checks}

The physical regressions sample finite poses, quaternions, velocities, and joint state;
bound wheel speed; check upright orientation and supported height; verify released-face
clearance; and hold the final snake beyond completion. Preflight rejects missing effort
support or undeclared limits before staging, and atomic command validation prevents a
partial actuator update.

These are necessary but not sufficient for robust control. They do not measure contact
penetration, energy, tire slip, solver residual, weld impulse, or sensitivity to
parameter perturbation.

\subsection{Determinism and reproducibility}

Semantic ordering is deterministic by construction. Floating-point physics is expected
to be reproducible within one pinned environment and scene, but no trace hash is
currently archived. Exact measured maxima and phase transition timestamps are also not
stored as an artifact. The regression protects thresholds and final state, which is
appropriate for development but insufficient for a comparative experimental paper.

Future result export should record at least:

\begin{itemize}
  \item simulation and wall timestamps;
  \item scenario phase/action/waypoint;
  \item module poses, joint targets/state/effort;
  \item connector position, axis, roll, and speed errors;
  \item contact count, force, penetration, and solver health;
  \item topology, event sequence, and constraint handles; and
  \item software version, dependency lock hash, scene configuration, and random seed.
\end{itemize}

\subsection{Test and CI analysis}

Coverage is branch-aware for the backend-neutral package. Optional \mujoco{} code is
checked separately because generated bindings lack complete stubs; Studio is checked by
its own strict configuration. CI installs only core dependencies in the Python matrix,
\mujoco{} in the headless job, and Studio in the Xvfb job. Package artifacts are built
and smoke-tested only after those jobs.

One current CI gap deserves explicit notice. The three tests in
\pathcode{tests/test_mujoco_contact_mechanics.py} validate authored tire pairs,
wheel-relative friction frames, and atomic weld/exclusion claim/release, but the
dedicated \mujoco{} job's explicit file list omits that file. They ran in the fresh local
537-test suite because \mujoco{} was installed; in a core-only CI job they skip. The file
should be added to the dedicated job before treating contact mechanics as continuously
protected on every pull request.

\subsection{Threats to validity}

\paragraph{Construct validity.}
Graph and event correctness are measured directly against the semantic contract.
Hardware accuracy is not: the contact model and actuators are provisional.

\paragraph{Internal validity.}
The root-motion guard establishes that physical scenarios do not teleport after
creation. Deterministic initial staging and authored routes nevertheless make the task
easier than arbitrary deployment.

\paragraph{External validity.}
Only one physical robot platform and one dynamic backend are exercised. The mock is a
semantic reference, not independent physics validation.

\paragraph{Statistical conclusion validity.}
The physical results are one tuned deterministic regression, not a repeated randomized
experiment. No confidence intervals or success-rate statistics can be claimed.

\paragraph{Artifact validity.}
The source checkout and lock file support reproduction, but there is no committed raw
trace, hardware calibration data set, public license, CITATION metadata, or explicit
public asset license. These are publication blockers for calling the artifact openly
reusable.

\subsection{Recommended benchmark campaign}

A meaningful next performance campaign should include:

\begin{enumerate}
  \item timestep sweep at 1, 2, and 5 ms;
  \item initial gap, lateral error, and yaw perturbation grid;
  \item repeated trace hashes and terminal metrics;
  \item requested versus achieved real-time factor with viewer off/on;
  \item scaling in modules/connectors and broad-phase density;
  \item model-view build/cache/serialization and IPC latency;
  \item contact-exclusion and friction-frame ablations; and
  \item calibrated hardware comparison for wheel and docking trajectories.
\end{enumerate}

